% ch06_numpy.tex — NumPy : Tableaux et Calcul Num\'erique
\chapter{NumPy --- Tableaux et Calcul Num\'erique}
\label{ch:numpy}

\begin{intuition}
Imaginez devoir mesurer la temp\'erature en 1\,000 points d'une plaque chauffante.
Avec une liste Python, chaque op\'eration n\'ecessite une boucle. NumPy\index{NumPy}
permet de traiter \emph{tous les points en une seule instruction}, comme un instrument
qui mesure toute la plaque d'un coup. C'est la diff\'erence entre compter des grains
de sable un par un et peser le tas entier.
\end{intuition}

% ============================================================
\section{Pourquoi NumPy~?}

\begin{codeblock}[title=\textbf{Python}]
\begin{minted}{python}
import time

# Comparaison : somme des carrés de 1 à 1 000 000
n = 1_000_000

# Avec une liste Python
debut = time.time()
resultat_liste = sum([i**2 for i in range(n)])
temps_liste = time.time() - debut

# Avec NumPy
import numpy as np
debut = time.time()
resultat_numpy = np.sum(np.arange(n)**2)
temps_numpy = time.time() - debut

print(f"Liste Python : {temps_liste:.4f} s")
print(f"NumPy        : {temps_numpy:.4f} s")
print(f"Accélération : ×{temps_liste/temps_numpy:.0f}")
\end{minted}
\end{codeblock}

\begin{codeoutput}
\begin{verbatim}
Liste Python : 0.2851 s
NumPy        : 0.0032 s
Accélération : ×89
\end{verbatim}
\end{codeoutput}

% ============================================================
\section{Cr\'eation de tableaux}
\index{np.array}\index{np.zeros}\index{np.ones}\index{np.linspace}\index{np.arange}

\begin{codeblock}[title=\textbf{Python}]
\begin{minted}{python}
import numpy as np

# À partir d'une liste
a = np.array([1.0, 2.0, 3.0, 4.0])
print(f"Array     : {a}")

# Tableaux spéciaux
zeros = np.zeros(5)
print(f"Zéros     : {zeros}")

uns = np.ones(4)
print(f"Uns       : {uns}")

# Séquence régulière
x = np.arange(0, 10, 2)
print(f"Arange    : {x}")

# N points équidistants
t = np.linspace(0, 1, 5)
print(f"Linspace  : {t}")
\end{minted}
\end{codeblock}

\begin{codeoutput}
\begin{verbatim}
Array     : [1. 2. 3. 4.]
Zéros     : [0. 0. 0. 0. 0.]
Uns       : [1. 1. 1. 1.]
Arange    : [0 2 4 6 8]
Linspace  : [0.   0.25 0.5  0.75 1.  ]
\end{verbatim}
\end{codeoutput}

% --- TikZ array indexing ---
\begin{figure}[ht]
\centering
\begin{tikzpicture}[
    cell/.style={rectangle, draw, minimum width=1.2cm, minimum height=0.8cm,
        fill=blue!10, align=center},
    idx/.style={font=\small\ttfamily, text=red!70!black},
]
\foreach \val [count=\i from 0] in {1.0, 2.0, 3.0, 4.0, 5.0} {
    \node[cell] (c\i) at (\i*1.3, 0) {\val};
    \node[idx, above=2pt of c\i] {\i};
    \pgfmathtruncatemacro{\neg}{\i - 5}
    \node[idx, below=2pt of c\i] {\neg};
}
\node[font=\small, left=0.5cm of c0, text=red!70!black] {indices};
\end{tikzpicture}
\caption{Indexation d'un tableau NumPy~: indices positifs (haut) et n\'egatifs (bas).}
\label{fig:numpy-indexation}
\end{figure}

% ============================================================
\section{Op\'erations \'el\'ement par \'el\'ement}
\index{op\'eration \'el\'ement par \'el\'ement}\index{broadcasting}

\begin{codeblock}[title=\textbf{Python}]
\begin{minted}{python}
import numpy as np

a = np.array([1, 2, 3, 4])
b = np.array([10, 20, 30, 40])

print(f"a + b   = {a + b}")
print(f"a * b   = {a * b}")
print(f"a ** 2  = {a ** 2}")
print(f"a + 100 = {a + 100}")    # broadcasting
print(f"a > 2   = {a > 2}")       # comparaison
\end{minted}
\end{codeblock}

\begin{codeoutput}
\begin{verbatim}
a + b   = [11 22 33 44]
a * b   = [ 10  40  90 160]
a ** 2  = [ 1  4  9 16]
a + 100 = [101 102 103 104]
a > 2   = [False False  True  True]
\end{verbatim}
\end{codeoutput}

\begin{definition}[Broadcasting]
Le \textbf{broadcasting}\index{broadcasting} permet \`a NumPy d'appliquer des op\'erations
entre tableaux de tailles diff\'erentes. Un scalaire est automatiquement
<<~\'etendu~>> \`a la taille du tableau.
\end{definition}

% ============================================================
\section{Indexation et slicing}

\begin{codeblock}[title=\textbf{Python}]
\begin{minted}{python}
import numpy as np

x = np.array([10, 20, 30, 40, 50, 60, 70])

print(f"Élément 3     : {x[3]}")
print(f"Slice [2:5]   : {x[2:5]}")

# Indexation booléenne (fancy indexing)
masque = x > 35
print(f"Masque        : {masque}")
print(f"Valeurs > 35  : {x[masque]}")

# Modification conditionnelle
x[x < 30] = 0
print(f"Après modif   : {x}")
\end{minted}
\end{codeblock}

\begin{codeoutput}
\begin{verbatim}
Élément 3     : 40
Slice [2:5]   : [30 40 50]
Masque        : [False False False  True  True  True  True]
Valeurs > 35  : [40 50 60 70]
Après modif   : [ 0  0 30 40 50 60 70]
\end{verbatim}
\end{codeoutput}

% ============================================================
\section{Fonctions math\'ematiques}
\index{np.sin}\index{np.cos}\index{np.exp}\index{np.log}\index{np.sqrt}

\begin{codeblock}[title=\textbf{Python}]
\begin{minted}{python}
import numpy as np

x = np.linspace(0, 2 * np.pi, 5)

print(f"x       = {np.round(x, 2)}")
print(f"sin(x)  = {np.round(np.sin(x), 4)}")
print(f"cos(x)  = {np.round(np.cos(x), 4)}")
print(f"exp(x)  = {np.round(np.exp(x), 2)}")

# Application : signal amorti
t = np.linspace(0, 5, 6)
signal = np.exp(-0.5 * t) * np.sin(2 * np.pi * t)
print(f"\nt       = {np.round(t, 1)}")
print(f"signal  = {np.round(signal, 4)}")
\end{minted}
\end{codeblock}

\begin{codeoutput}
\begin{verbatim}
x       = [0.   1.57 3.14 4.71 6.28]
sin(x)  = [ 0.      1.      0.     -1.     -0.    ]
cos(x)  = [ 1.      0.     -1.     -0.      1.    ]
exp(x)  = [  1.     4.81  23.14 111.32 535.49]

t       = [0. 1. 2. 3. 4. 5.]
signal  = [ 0.     -0.      0.     -0.      0.     -0.    ]
\end{verbatim}
\end{codeoutput}

% ============================================================
\section{Statistiques descriptives}
\index{np.mean}\index{np.std}\index{np.min}\index{np.max}

\begin{codeblock}[title=\textbf{Python}]
\begin{minted}{python}
import numpy as np

# Mesures expérimentales (concentration en mg/L)
mesures = np.array([12.3, 11.8, 12.5, 11.9, 12.1, 12.4, 12.0, 11.7])

print(f"Moyenne     : {np.mean(mesures):.2f} mg/L")
print(f"Écart-type  : {np.std(mesures):.2f} mg/L")
print(f"Minimum     : {np.min(mesures):.2f} mg/L")
print(f"Maximum     : {np.max(mesures):.2f} mg/L")
print(f"Médiane     : {np.median(mesures):.2f} mg/L")
\end{minted}
\end{codeblock}

\begin{codeoutput}
\begin{verbatim}
Moyenne     : 12.09 mg/L
Écart-type  : 0.27 mg/L
Minimum     : 11.70 mg/L
Maximum     : 12.50 mg/L
Médiane     : 12.05 mg/L
\end{verbatim}
\end{codeoutput}

% ============================================================
\section{Alg\`ebre lin\'eaire \'el\'ementaire}
\index{np.dot}\index{multiplication matricielle}

\begin{codeblock}[title=\textbf{Python}]
\begin{minted}{python}
import numpy as np

# Matrices 2×2
A = np.array([[1, 2],
              [3, 4]])

B = np.array([[5, 6],
              [7, 8]])

# Produit matriciel
C = A @ B    # équivalent à np.dot(A, B)
print("A @ B =")
print(C)

# Déterminant et inverse
det_A = np.linalg.det(A)
inv_A = np.linalg.inv(A)
print(f"\ndet(A) = {det_A:.1f}")
print(f"A⁻¹ =\n{inv_A}")
\end{minted}
\end{codeblock}

\begin{codeoutput}
\begin{verbatim}
A @ B =
[[19 22]
 [43 50]]

det(A) = -2.0
A⁻¹ =
[[-2.   1. ]
 [ 1.5 -0.5]]
\end{verbatim}
\end{codeoutput}

% ============================================================
\section{Nombres al\'eatoires}
\index{np.random}

\begin{codeblock}[title=\textbf{Python}]
\begin{minted}{python}
import numpy as np

rng = np.random.default_rng(seed=42)

# Uniforme entre 0 et 1
u = rng.uniform(0, 1, size=5)
print(f"Uniforme  : {np.round(u, 3)}")

# Distribution normale (moyenne=0, écart-type=1)
n = rng.normal(0, 1, size=5)
print(f"Normale   : {np.round(n, 3)}")

# Entiers aléatoires (lancer de dé)
des = rng.integers(1, 7, size=10)
print(f"Lancers   : {des}")
\end{minted}
\end{codeblock}

\begin{codeoutput}
\begin{verbatim}
Uniforme  : [0.773 0.438 0.859 0.697 0.094]
Normale   : [ 0.314  0.459 -0.537  0.268 -0.232]
Lancers   : [5 2 3 6 1 4 2 5 3 6]
\end{verbatim}
\end{codeoutput}

% ============================================================
\section{Erreurs courantes}

\begin{errmark}
\begin{enumerate}
    \item \textbf{Incompatibilit\'e de dimensions (\py{shape mismatch})}~: op\'erer sur
    des tableaux de tailles incompatibles provoque une \py{ValueError}.
    \begin{minted}{python}
a = np.array([1, 2, 3])
b = np.array([1, 2])
# a + b → ValueError: operands could not be broadcast
    \end{minted}

    \item \textbf{Oublier les op\'erations \'el\'ement par \'el\'ement}~: l'op\'erateur \py{*}
    entre deux tableaux fait une multiplication \'el\'ement par \'el\'ement, \emph{pas}
    un produit matriciel. Utilisez \py{@} ou \py{np.dot} pour le produit matriciel.

    \item \textbf{Confusion \py{np.arange} / \py{np.linspace}}~:
    \py{np.arange(0, 1, 0.1)} sp\'ecifie un \emph{pas},
    \py{np.linspace(0, 1, 10)} sp\'ecifie un \emph{nombre de points}.

    \item \textbf{Modifier une vue au lieu d'une copie}~: le slicing NumPy cr\'ee une
    \emph{vue} (pas une copie). Modifier la vue modifie l'original~!
    \begin{minted}{python}
a = np.array([1, 2, 3, 4])
b = a[1:3]     # vue, pas copie
b[0] = 99      # a devient [1, 99, 3, 4] !
b = a[1:3].copy()  # copie indépendante
    \end{minted}
\end{enumerate}
\end{errmark}

% ============================================================
\section{Mini-projet~: Estimation de $\pi$ par Monte Carlo}

\begin{projet}
La m\'ethode de Monte Carlo\index{Monte Carlo} estime $\pi$ en tirant des points
al\'eatoires dans un carr\'e $[-1, 1]^2$ et en comptant ceux qui tombent dans le
cercle unit\'e ($x^2 + y^2 \leq 1$). Le rapport donne~:
\[
    \pi \approx 4 \times \frac{\text{points dans le cercle}}{\text{points totaux}}
\]
\end{projet}

\begin{codeblock}[title=\textbf{Python}]
\begin{minted}{python}
import numpy as np

rng = np.random.default_rng(seed=42)
N = 1_000_000

# Tirer N points aléatoires dans [-1, 1] × [-1, 1]
x = rng.uniform(-1, 1, size=N)
y = rng.uniform(-1, 1, size=N)

# Compter les points dans le cercle unité
dans_cercle = x**2 + y**2 <= 1
nb_dans_cercle = np.sum(dans_cercle)

# Estimation de pi
pi_estime = 4 * nb_dans_cercle / N
erreur = abs(pi_estime - np.pi)

print(f"Points totaux     : {N:,}")
print(f"Points dans cercle: {nb_dans_cercle:,}")
print(f"π estimé          : {pi_estime:.6f}")
print(f"π exact           : {np.pi:.6f}")
print(f"Erreur            : {erreur:.6f}")
\end{minted}
\end{codeblock}

\begin{codeoutput}
\begin{verbatim}
Points totaux     : 1,000,000
Points dans cercle: 785,436
π estimé          : 3.141744
π exact           : 3.141593
Erreur            : 0.000151
\end{verbatim}
\end{codeoutput}

% ============================================================
\section{Exercices}

\begin{exercice}[$\star$ --- Cr\'eation de tableaux]
Cr\'eez les tableaux suivants~: (a) les entiers de 0 \`a 99, (b) 50 valeurs
\'equidistantes entre $-\pi$ et $\pi$, (c) une matrice $3\times 3$ de z\'eros
avec des 1 sur la diagonale (matrice identit\'e).
\end{exercice}

\begin{exercice}[$\star$ --- Conversion de temp\'eratures]
Cr\'eez un tableau de temp\'eratures en Celsius de $-40$ \`a $100$ (pas de 10).
Convertissez-le en Fahrenheit ($F = 1{,}8 \times C + 32$) \emph{sans boucle}.
\end{exercice}

\begin{exercice}[$\star\star$ --- Analyse de donn\'ees exp\'erimentales]
G\'en\'erez 1000 mesures simul\'ees d'une distribution normale ($\mu = 50$, $\sigma = 5$).
Calculez la moyenne, l'\'ecart-type, et comptez combien de mesures sont \`a plus de
$2\sigma$ de la moyenne. Comparez avec la pr\'ediction th\'eorique ($\approx 5\%$).
\end{exercice}

\begin{exercice}[$\star\star$ --- Produit scalaire]
Calculez l'angle entre les vecteurs $\vec{u} = (1, 2, 3)$ et $\vec{v} = (4, 5, 6)$
en utilisant la formule $\cos\theta = \frac{\vec{u}\cdot\vec{v}}{|\vec{u}||\vec{v}|}$.
\end{exercice}

\begin{exercice}[$\star\star\star$ --- R\'esolution de syst\`eme lin\'eaire]
R\'esolvez le syst\`eme d'\'equations lin\'eaires suivant avec \py{np.linalg.solve}~:
\[
\begin{cases}
2x + 3y - z = 1 \\
4x + y + 2z = 2 \\
-x + 2y + 3z = 3
\end{cases}
\]
V\'erifiez votre r\'esultat en calculant $A\vec{x}$ et en comparant avec $\vec{b}$.
\end{exercice}

% ============================================================
\begin{formules}{R\'esum\'e du chapitre~: NumPy}
\begin{itemize}
    \item \textbf{Cr\'eation}~: \py{np.array}, \py{np.zeros}, \py{np.ones},
    \py{np.linspace}, \py{np.arange}
    \item \textbf{Op\'erations}~: \py{+}, \py{-}, \py{*}, \py{/}, \py{**}
    s'appliquent \'el\'ement par \'el\'ement
    \item \textbf{Broadcasting}~: un scalaire est \'etendu automatiquement
    \item \textbf{Maths}~: \py{np.sin}, \py{np.cos}, \py{np.exp}, \py{np.log}, \py{np.sqrt}
    \item \textbf{Stats}~: \py{np.mean}, \py{np.std}, \py{np.min}, \py{np.max}, \py{np.median}
    \item \textbf{Alg\`ebre}~: \py{A @ B} (produit matriciel), \py{np.linalg.det},
    \py{np.linalg.inv}, \py{np.linalg.solve}
    \item \textbf{Al\'eatoire}~: \py{rng = np.random.default\_rng(seed)}
    \item \textbf{Indexation bool\'eenne}~: \py{x[x > 0]} filtre les \'el\'ements positifs
\end{itemize}
\end{formules}
